% Section 5.2, from lectures/L05/appendix/b_type_traits_and_constraints.md.

\section{Type Traits and Constraints}
\label{c5:sec:traits}\label{c5:app:b}
Type traits allow programs to inspect and reason about types at compile time.

They are provided in the C++ standard library via the header:

\begin{cppcode}
#include <type_traits>
\end{cppcode}

\noindent Type traits are commonly used together with templates to create safer and more robust
code, for instance to prevent misuse of low-level hardware APIs (e.g.\ restricting register
operations to integral types).

\subsection{Common type traits}
\label{c5:sec:commontraits}
Some commonly used type traits include:
\begin{itemize}
  \item \code{std::is_integral} --- checks if a type is an integer type.
  \item \code{std::is_arithmetic} --- checks if a type is numeric (integer or floating point).
  \item \code{std::is_unsigned} --- checks if a type is an unsigned integer.
  \item \code{std::is_floating_point} --- checks if a type is a floating-point type.
\end{itemize}
Each type trait provides a boolean value:

\begin{cppcode}
std::is_integral<int>::value     // true
std::is_integral<double>::value  // false
\end{cppcode}

\noindent Since C++17, a shorthand form is also available:

\begin{cppcode}
std::is_integral_v<int> // true
\end{cppcode}

\subsection{Compile-time constraints with \code{static_assert()}}
\label{c5:sec:staticassert}
Type traits are often used together with \code{static_assert()} to enforce constraints at compile
time.

Example:

\begin{cppcode}
#include <type_traits>

/**
 * @brief Add two numbers.
 *
 * @tparam T The numeric type. Must be arithmetic.
 */
template<typename T>
constexpr T add(const T x, const T y) noexcept
{
    static_assert(std::is_arithmetic<T>::value, "Cannot add values of non-arithmetic types!");
    return x + y;
}
\end{cppcode}

\noindent If the function is used incorrectly:

\begin{cppcode}
add("Hello", "World");
\end{cppcode}

\noindent The \code{static_assert()} fails and the compiler stops with the error message given
above. Without it, the compiler would still reject the call, since two pointers cannot be added,
but with a less readable error from inside the function body.

\subsection{Why compile-time checks matter in embedded systems}
\label{c5:sec:whycompiletime}
In embedded systems:
\begin{itemize}
  \item Errors should be detected as early as possible.
  \item Debugging run-time issues can be difficult.
  \item Safety and robustness are critical.
\end{itemize}
Compile-time checks provide:
\begin{itemize}
  \item Early error detection.
  \item Clear diagnostics.
  \item Zero run-time overhead.
\end{itemize}

\subsection{Creating custom type traits}
\label{c5:sec:customtraits}
If no suitable type trait exists, we can create our own.

A type trait is typically implemented as a template struct that exposes a boolean constant named
\code{value}. We set the constant \code{value} to \code{false} by default for a given type
\code{T}:

\begin{cppcode}
template<typename T>
struct isUnsigned
{
    static constexpr bool value{false};
};
\end{cppcode}

\noindent We then set the constant \code{value} to \code{true} for unsigned types, such as
\code{std::uint8_t} and \code{std::uint32_t}:

\begin{cppcode}
#include <cstdint>

template<>
struct isUnsigned<std::uint8_t>
{
    static constexpr bool value{true};
};

template<>
struct isUnsigned<std::uint16_t>
{
    static constexpr bool value{true};
};

template<>
struct isUnsigned<std::uint32_t>
{
    static constexpr bool value{true};
};
\end{cppcode}

\noindent Usage:

\begin{cppcode}
// Generate a compiler error if T isn't of unsigned type:
static_assert(isUnsigned<T>::value, "T must be unsigned!");
\end{cppcode}

\subsection{Variable templates (C++17)}
\label{c5:sec:variabletemplates}
For convenience, we can define a variable template:

\begin{cppcode}
template<typename T>
inline constexpr bool isUnsigned_v{isUnsigned<T>::value};
\end{cppcode}

\noindent This allows:

\begin{cppcode}
// Generate a compiler error if T isn't of unsigned type:
static_assert(isUnsigned_v<T>, "T must be unsigned!");
\end{cppcode}

\note In practice, the standard library already provides many useful traits (e.g.\
\code{std::is_unsigned}). Custom traits are mainly used when no suitable standard trait exists.

\subsection{Summary}
\label{c5:sec:traitssummary}
\begin{itemize}
  \item Type traits allow inspection of types at compile time.
  \item \code{static_assert()} enables compile-time validation.
  \item Compile-time checks improve safety and robustness.
  \item Custom type traits can be created when needed.
  \item These tools are essential for building safe template-based APIs.
\end{itemize}
